\documentclass{article}
\usepackage[T1]{fontenc}
\usepackage{lmodern}
\usepackage{graphicx}
\usepackage{amsmath,amssymb}
\usepackage[a4paper,margin=2cm]{geometry}
\usepackage[absolute,overlay]{textpos}
\setlength{\TPHorizModule}{1mm}
\setlength{\TPVertModule}{1mm}
\usepackage[indonesian]{babel}
\usepackage{hyperref}
\setlength{\emergencystretch}{2em}
% unit: Halaman 9

\begin{document}

% === Halaman 9 (python/native) ===

Rinaldi Munir/II4021 Kriptografi

9

Medan (Field)

• Medan (field) adalah system aljabar yang terdidi dari


- himpunan elemen (disimbolkan dengan F) 


- operasi biner, yaitu penjumlahan (+) dan perkalian ().

• Sebuah struktur aljabar <F, +, > disebut medan jika dan hanya jika:

 1. <F, +> adalah grup abelian

 2. <F – \{0\}, > adalah grup abelian

 3. Operasi  menyebar terhadap operasi +  (sifat distributif)

     Distributif: x  ( y + z) = (x  y) + (x  z)



(x + y)  z  = (x  z) + (y  z)

• Jadi, sebuah medan memenuhi aksioma: closure, komutatif, asosiatif, dan 

distributif

\end{document}
